% !PRINTED 126
\weberpage{126}
%<<CONT

different from zero. Thus suppose that $f'(\alpha),f''(\alpha),\ldots,f^{(m-1)}(\alpha)$ vanish, but $f^{(m)}(\alpha)$ does not vanish. We then have, by \S\,13,
\begin{equation}
 f(\alpha+h)=f(\alpha)+\frac{h^m}{\Pi(m)}f^{(m)}(\alpha)
 +\frac{h^{m+1}}{\Pi(m+1)}f^{(m+1)}(\alpha)+\cdots .
\tag{12}
\end{equation}

We choose, as is always possible by \S\,31, V, a positive number $\varepsilon$ so that
\begin{equation}
\left|\frac{h}{m+1}\frac{f^{(m+1)}(\alpha)}{f^{(m)}(\alpha)}
+\frac{h^2}{(m+1)(m+2)}\frac{f^{(m+2)}(\alpha)}{f^{(m)}(\alpha)}+\cdots\right|<1
\tag{13}
\end{equation}
whenever
\begin{equation}
 |h|<\varepsilon,
\tag{14}
\end{equation}
and a positive proper fraction $\delta$ such that
\begin{equation}
 \delta |f(\alpha)|<\left|\frac{\varepsilon^m f^{(m)}(\alpha)}{\Pi(m)}\right|,
\tag{15}
\end{equation}
which, if $f(\alpha)$ does not vanish, is likewise possible. We then determine $h$ (on the basis of \S\,32) from the equation
\begin{equation}
 \frac{h^m f^{(m)}(\alpha)}{\Pi(m)}=-\delta f(\alpha),
\tag{16}
\end{equation}
whence, by (15),
\[
 |h|<\varepsilon
\]
follows. From (12), when the value from (16) is substituted for $h^m$, we now get
\[
 f(\alpha+h)=(1-\delta)f(\alpha)
 -\delta f(\alpha)\left(\frac{h}{m+1}\frac{f^{(m+1)}(\alpha)}{f^{(m)}(\alpha)}
 +\frac{h^2}{(m+1)(m+2)}\frac{f^{(m+2)}(\alpha)}{f^{(m)}(\alpha)}+\cdots\right),
\]
and hence
\begin{align}
 |f(\alpha+h)|\le{}& |f(\alpha)|(1-\delta)\notag\\
 &+\delta |f(\alpha)|\left|\frac{h}{m+1}\frac{f^{(m+1)}(\alpha)}{f^{(m)}(\alpha)}
 +\frac{h^2}{(m+1)(m+2)}\frac{f^{(m+2)}(\alpha)}{f^{(m)}(\alpha)}+\cdots\right|.
\tag{17}
\end{align}
By (13), therefore,
\begin{equation}
 |f(\alpha+h)|<|f(\alpha)|.\tag{18}
\end{equation}
\noindent\hfill q.e.d.

Thus the proof of the fundamental theorem is completed.
